% =============================================================
%  HEADER — Dossier_Final_groupe17 |  ENSAI 2026-2027
%  Projet : Laborscope
%  Les "chapitres" sont des \section redéfinies visuellement.
% =============================================================

% ---------- Langue ----------
%\usepackage[french]{babel}
\addto\captionsfrench{%
  \renewcommand{\listfigurename}{Liste des figures}%
  \renewcommand{\listtablename}{Liste des tableaux}%
  \renewcommand{\contentsname}{Table des matières}%
  \renewcommand{\figurename}{Figure}%
  \renewcommand{\tablename}{Tableau}%
}
\AtBeginDocument{%
  \renewcommand{\listfigurename}{Liste des figures}%
  \renewcommand{\listtablename}{Liste des tableaux}%
  \renewcommand{\contentsname}{Table des matières}%
  \renewcommand{\figurename}{Figure}%
  \renewcommand{\tablename}{Tableau}%
  \hypersetup{linkcolor=black}%
}

% ---------- Polices ----------
\usepackage{fontspec}
\usepackage{unicode-math}
%\usepackage{lmodern}
%\setmainfont{Source Sans Pro}

% ---------- Interligne ----------
\usepackage{setspace}
\onehalfspacing


% ---------- Maths ----------
\usepackage{amsmath}
\usepackage{amssymb}
\everymath{\displaystyle}

% ---------- Graphiques & tableaux ----------
\usepackage{graphicx}
\usepackage{float}
\usepackage{array}
\usepackage{booktabs}
\usepackage{multirow}
\usepackage{longtable}
\usepackage{threeparttable}
\renewcommand{\arraystretch}{1.5}
\setlength{\textfloatsep}{8pt}
\setlength{\intextsep}{8pt}

% ---------- PDF ----------
\usepackage{pdfpages}

% =============================================================
%  PALETTE
% =============================================================
\usepackage{xcolor}
\definecolor{MainColor}{HTML}{1565C0}      % bleu principal
\definecolor{AccentColor}{HTML}{7B2FBE}    % violet du cadre titre
\definecolor{MainLight}{HTML}{E0F4FF}      % bleu très clair fond
\definecolor{NeutralColor}{HTML}{546E7A}   % gris-bleu neutre
\definecolor{AccentLight}{HTML}{F3E8FF}    % violet très clair
\definecolor{PurpleColor}{HTML}{FF6F00}    % orange accent
\definecolor{PurpleLight}{HTML}{FFF3E0}    % orange très clair
\definecolor{bg}{rgb}{0.96,0.96,0.96}

% ---------- Tableaux ----------
\usepackage{colortbl}
\arrayrulecolor{white}

% ---------- Minted (un seul chargement) ----------
\usepackage{minted}
\makeatletter\let\listoflistings\@undefined\makeatother
\usemintedstyle{friendly}
\setminted{
  bgcolor=bg, fontsize=\small, linenos=true,
  numbersep=6pt, framesep=2mm, breaklines=true, baselinestretch=1.1
}
\renewcommand{\theFancyVerbLine}{%
  \textcolor{NeutralColor}{\small\arabic{FancyVerbLine}}%
}

% ---------- Hyperref ----------
\usepackage[colorlinks=true, linkcolor=black,
            urlcolor=AccentColor, citecolor=MainColor]{hyperref}
\usepackage[french]{cleveref}

% ---------- Mise en page ----------
\usepackage{geometry}
\geometry{a4paper, left=25mm, right=25mm, top=25mm, bottom=25mm}

% ---------- TikZ & tcolorbox ----------
\usepackage{tikz}
\usetikzlibrary{shapes.geometric}
\usepackage[most]{tcolorbox}

\tcbset{
  isebase/.style={
    coltitle=white, fonttitle=\bfseries\small,
    arc=3pt, boxrule=0pt,
    left=8pt, right=8pt, top=5pt, bottom=5pt, titlerule=0pt,
    attach boxed title to top left={yshift=-2mm, xshift=6pt},
    boxed title style={arc=2pt, boxrule=0pt, left=4pt, right=4pt},
    breakable
  }
}

\newtcolorbox{definition}[1][]{isebase,enhanced,
  title={Définition}, colback=MainLight, colframe=MainColor,
  colbacktitle=MainColor, frame hidden,
  borderline west={3pt}{0pt}{MainColor}, #1}

\newtcolorbox{theoreme}[1][]{isebase,enhanced,
  title={Théorème}, colback=MainLight, colframe=MainColor,
  colbacktitle=MainColor, frame hidden,
  borderline west={3pt}{0pt}{MainColor}, #1}

\newtcolorbox{modele}[1][]{isebase,enhanced,
  title={Modèle}, colback=MainLight, colframe=MainColor,
  colbacktitle=MainColor, frame hidden,
  borderline west={3pt}{0pt}{MainColor}, #1}

\newtcolorbox{hypotheses}[1][]{isebase,enhanced,
  title={Hypothèses}, colback=MainLight, colframe=MainColor,
  colbacktitle=MainColor, frame hidden,
  borderline west={3pt}{0pt}{MainColor}, #1}

% decision = même style que hypotheses
\newtcolorbox{decision}[1][]{isebase,enhanced,
  title={Règle de décision}, colback=MainLight, colframe=MainColor,
  colbacktitle=MainColor, frame hidden,
  borderline west={3pt}{0pt}{MainColor}, #1}

\newtcolorbox{remarque}[1][]{isebase,enhanced,
  title={Remarque}, colback=AccentLight, colframe=AccentColor,
  colbacktitle=AccentColor, frame hidden,
  borderline west={3pt}{0pt}{AccentColor}, #1}

\newtcolorbox{attention}[1][]{isebase,enhanced,
  title={Attention}, colback=AccentLight, colframe=AccentColor,
  colbacktitle=AccentColor, frame hidden,
  borderline west={4pt}{0pt}{AccentColor}, #1}

\newtcolorbox{resultat}[1][]{isebase,enhanced,
  title={Résultat}, colback=PurpleLight, colframe=PurpleColor,
  colbacktitle=PurpleColor, frame hidden,
  borderline west={3pt}{0pt}{PurpleColor}, #1}

% ---------- En-têtes & pieds ----------
\usepackage{fancyhdr}
\pagestyle{fancy}
\fancyhf{}
\renewcommand{\headrulewidth}{0pt}
\renewcommand{\footrulewidth}{0pt}
%\renewcommand{\headrule}{%
%  \vspace{-4pt}{\color{MainColor}\rule{\headwidth}{1.5pt}}}

\fancyhead[L]{}
%\fancyhead[C]{\small\color{NeutralColor}Déterminants macroéconomiques de la croissance du PIB de la Sierra Leone}
\fancyhead[R]{}
\fancyfoot[C]{%
  \begin{tikzpicture}[baseline=-0.5ex]
    \node[circle, draw=MainColor, fill=white, text=MainColor,
          inner sep=2pt, font=\small\bfseries, line width=0.8pt]{\thepage};
  \end{tikzpicture}}
\fancyfoot[L]{}
\fancyfoot[R]{\small\color{NeutralColor}\textcopyright\ ENSAI 2026-2027}

\fancypagestyle{plain}{%
  \fancyhf{}
  \renewcommand{\headrulewidth}{0pt}
  \renewcommand{\footrulewidth}{0pt}
  \fancyfoot[L]{%
    \begin{tikzpicture}[baseline=-0.5ex]
      \node[circle, draw=MainColor, fill=white, text=MainColor,
            inner sep=2pt, font=\small\bfseries, line width=0.8pt]{\thepage};
    \end{tikzpicture}}
  \fancyfoot[C]{\small\color{NeutralColor}\textcopyright\ ENSAI 2026-2027}
}

% =============================================================
%  FORMATAGE DES TITRES — titlesec + compteurs indépendants
%  Rapport : Déterminants macroéconomiques du PIB — Sierra Leone
%
%  CAUSE DU BUG "1.0.0" :
%  \titleformat{\section}[block] passe le titre à \sectiontitre
%  AVANT d'incrémenter \arabic{section}. Donc \thesection affiche
%  \chapnum.0 au lieu de \chapnum.1 pour la première section.
%
%  SOLUTION : on n'utilise pas \arabic{section} du tout.
%  On crée des compteurs indépendants (isesecnum, isesubsecnum)
%  qu'on incrémente NOUS-MÊMES au début de \sectiontitre /
%  \subsectiontitre — après le rendu réel de \mychapter.
%
%  SUPPRESSION DES DOUBLONS TDC :
%  scrartcl ajoute automatiquement des entrées TDC pour chaque
%  \section/\subsection/\subsubsection. Nos commandes ajoutent
%  déjà les entrées manuellement via \addcontentsline.
%  On désactive les entrées automatiques en mettant tocdepth à -2.
%
%  STRUCTURE DU RAPPORT :
%  Introduction → Chapitre 1 (Théorie & Littérature) →
%  Chapitre 2 (Méthodologie) → Chapitre 3 (Résultats) →
%  Conclusion → Références → Annexe (Code R)
% =============================================================
\usepackage{titlesec}

% Désactiver les entrées TDC automatiques de scrartcl
% (nos \addcontentsline manuels dans \sectiontitre etc. suffisent)
\setcounter{tocdepth}{-2}
\setcounter{secnumdepth}{-2}

% ── Compteur chapitre ─────────────────────────────────────────
\newcounter{chapnum}
\renewcommand{\thechapnum}{\arabic{chapnum}}

% ── Compteurs sections (remis à 0 par \mychapter / \sectiontitre)
\newcounter{isesecnum}
\renewcommand{\theisesecnum}{\thechapnum.\arabic{isesecnum}}

\newcounter{isesubsecnum}
\renewcommand{\theisesubsecnum}{\theisesecnum.\arabic{isesubsecnum}}

\newcounter{isesubsubsecnum}
\renewcommand{\theisesubsubsecnum}{\theisesubsecnum.\arabic{isesubsubsecnum}}

% ── \mychapter : badge numéroté ───────────────────────────────
\newcommand{\mychapter}[1]{%
  \refstepcounter{chapnum}%
  \setcounter{isesecnum}{0}%
  \setcounter{isesubsecnum}{0}%
  \setcounter{isesubsubsecnum}{0}%
  \addcontentsline{toc}{chapter}{\protect\numberline{\thechapnum}#1}%
  \vspace{30pt}%
  \noindent
  \begin{tikzpicture}
    \node[text=MainColor, font=\small\itshape, anchor=south west]
      at (0,4.0em) {Chapitre};
    \node[fill=MainColor, text=white,
          font=\fontsize{48}{48}\selectfont\bfseries,
          minimum width=1.6em, minimum height=1.6em,
          inner sep=6pt, anchor=south west] (nb) at (0,0) {\thechapnum};
    \node[text=MainColor, font=\Large\bfseries,
          anchor=south west, text width=0.72\textwidth, align=left]
      at ([xshift=12pt] nb.south east) {\MakeUppercase{#1}};
  \end{tikzpicture}%
  \vspace{6pt}\par
}

% ── \mychapterstar : chapitre non numéroté ────────────────────
\newcommand{\mychapterstar}[1]{%
  \addcontentsline{toc}{chapter}{#1}%
  \vspace{4pt}%
  \noindent{\Large\bfseries\color{MainColor}\MakeUppercase{#1}}\par
  \vspace{4pt}\par%
}

% ── Sections (## dans Quarto → \subsection dans scrartcl) ────
% Dans scrartcl : # → \section, ## → \subsection, ### → \subsubsection
% Nos "chapitres" sont gérés manuellement via \mychapter.
% Donc ## correspond visuellement à nos sections numérotées.
\newcommand{\sectiontitre}[1]{%
  \refstepcounter{isesecnum}%
  \setcounter{isesubsecnum}{0}%
  \setcounter{isesubsubsecnum}{0}%
  \addcontentsline{toc}{section}{\theisesecnum\quad #1}%
  \noindent
  \begin{tikzpicture}
    \node[fill=MainColor, text=white, font=\bfseries\normalsize,
          inner sep=5pt, minimum width=2.4em, minimum height=1.6em,
          anchor=base west] (sn) at (0,0) {\theisesecnum};
    \node[text=MainColor, font=\bfseries\normalsize, anchor=base west]
      at ([xshift=8pt] sn.base east) {\MakeUppercase{#1}};
  \end{tikzpicture}%
  \vspace{-2pt}{\color{MainColor}\hrule height 0.8pt}\vspace{2pt}%
}
\titleformat{\subsection}[block]{}{}{0pt}{\sectiontitre}
\titlespacing*{\subsection}{0pt}{14pt}{2pt}

% ── Sous-sections (### dans Quarto → \subsubsection dans scrartcl) ────
\newcommand{\subsectiontitre}[1]{%
  \refstepcounter{isesubsecnum}%
  \setcounter{isesubsubsecnum}{0}%
  \addcontentsline{toc}{subsection}{\theisesubsecnum\quad #1}%
  \noindent
  \begin{tikzpicture}
    \node[fill=MainColor, text=white, font=\bfseries\small,
          inner sep=4pt, minimum width=3em, minimum height=1.4em,
          anchor=base west] (ss) at (0,0) {\theisesubsecnum};
    \node[text=MainColor, font=\bfseries\small, anchor=base west]
      at ([xshift=6pt] ss.base east) {#1};
  \end{tikzpicture}%
  \vspace{-2pt}{\color{MainColor!40}\hrule height 0.6pt}\vspace{2pt}%
}
\titleformat{\subsubsection}[block]{}{}{0pt}{\subsectiontitre}
\titlespacing*{\subsubsection}{0pt}{5pt}{2pt}

% ── Sous-sous-sections (#### dans Quarto → \paragraph dans scrartcl) ──
\newcommand{\subsubsectiontitre}[1]{%
  \refstepcounter{isesubsubsecnum}%
  \noindent{\normalsize\itshape\color{NeutralColor}\theisesubsubsecnum\enspace #1}%
}
\titleformat{\paragraph}[block]{}{}{0pt}{\subsubsectiontitre}
\titlespacing*{\paragraph}{0pt}{6pt}{2pt}

% -------------------------------------------------------
% Entrées TDM pour le niveau "chapter" manuel
% -------------------------------------------------------
%  TDM — chapter/section/subsection tous en mode manuel
%  (les entrées viennent de \addcontentsline dans les
%   commandes \mychapter, \sectiontitre, \subsectiontitre)
% -------------------------------------------------------
\usepackage{tocloft}

\renewcommand{\cfttoctitlefont}{\Huge\bfseries\color{black}}
\renewcommand{\cftloftitlefont}{\Huge\bfseries\color{black}}
\renewcommand{\cftlottitlefont}{\Huge\bfseries\color{black}}
\renewcommand{\cftdotsep}{3}

\makeatletter
% Style pour les chapitres (ajoutés via \addcontentsline{toc}{chapter})
\def\l@chapter#1#2{%
  \addpenalty{-\@highpenalty}%
  \vskip 8pt \@plus\p@
  \begingroup
    \parindent\z@ \rightskip\@pnumwidth
    \parfillskip-\@pnumwidth
    {\bfseries\color{black} #1}\nobreak\hfill\nobreak
    \hb@xt@\@pnumwidth{\hss\bfseries\color{black}#2}\par
    \penalty\@highpenalty
  \endgroup
}
% Style pour les sections (ajoutées via \addcontentsline{toc}{section})
\def\l@section#1#2{%
  \vskip 3pt
  \begingroup
    \parindent\z@ \rightskip\@pnumwidth
    \parfillskip-\@pnumwidth
    \hspace{1.5em}{\color{black} #1}\nobreak\dotfill\nobreak
    \hb@xt@\@pnumwidth{\hss\color{black}#2}\par
  \endgroup
}
% Style pour les sous-sections (ajoutées via \addcontentsline{toc}{subsection})
\def\l@subsection#1#2{%
  \vskip 2pt
  \begingroup
    \parindent\z@ \rightskip\@pnumwidth
    \parfillskip-\@pnumwidth
    \hspace{3em}{\small\color{black} #1}\nobreak\dotfill\nobreak
    \hb@xt@\@pnumwidth{\hss\small\color{black}#2}\par
  \endgroup
}
\makeatother

% ---------- Listes ----------
\usepackage{enumitem,pifont}
\setlist[itemize,1]{label={\color{AccentColor}\textbullet}}
\setlist[itemize,2]{label={\color{MainColor}$\circ$}}

\AtBeginDocument{
  \floatplacement{table}{H}
  \floatplacement{figure}{H}
}